\section{A theory of prefix-cache metastability}
\label{sec:theory}

This section derives the failure structure that Sections \ref{sec:collapse}--\ref{sec:span}
measure: a retention fixed point coupling hit rate to write rate
(\ref{sec:retention}), the occupancy-and-queueing feedback that creates the cold
branch under this workload's reuse gaps (\ref{sec:occupancy}), first-order
break-point and tier-sizing estimates, each checked against a
measured quantity (\ref{sec:breakpoints}--\ref{sec:tier-theory}), and the dimensionless groups that
organize the regimes (\ref{sec:dimensionless}). The theory is deliberately first-order:
it predicts the existence, location, and scaling of the regimes; the
calibrated simulation of Section~\ref{sec:models} carries the quantitative claims.

\paragraph{Notation.} $N$ replicas each hold an HBM KV pool of $C$ tokens under LRU
block eviction (fleet pool $S = N \cdot C$) and optionally a host-memory
tier of $\Ccpu$ tokens per replica; $p$ is the per-replica prefill rate
and $\Br$ the effective per-replica restore rate (tokens/s). Sessions
arrive at rate $\ls$ and issue $\Ereq$ requests over their
lifetime, so the long-run request rate is $\lr = \ls \cdot \Ereq$.
A request with context $c$ prefills its new input $\delta$ on a
hit and $c + \delta$ on a miss, and appends its output to the cache
either way. $F(t)$ is the reuse-gap distribution: the probability that
a session's next request arrives within $t$ seconds of its predecessor
becoming reusable. Calibrated values for the measured stack ($N = 8$,
$C = 1.66$M tokens, $S = 13.28$M, $p = 17$k tokens/s, $\Ereq = 118$, reuse
gaps capped at 10.5\,s by the protocol): mean write per miss
$\wmiss = 117.8$k tokens, mean write per hit $\whit = 5.2$k, mean
prefill per miss $\cmiss = 116.9$k (workload functionals derived from
the trace-calibrated spec; Section~\ref{sec:calibration}).

\subsection{The retention fixed point}
\label{sec:retention}

An LRU cache under aggregate write rate $W$ retains an entry for the
characteristic time $T = S / W$: new writes traverse the pool in $T$
seconds, and an entry survives if touched again within that window
\cite{che2002hierarchical,fricker2012versatile}. In a prefix cache, $W$ is
not exogenous. At hit rate $h$,
\begin{equation}
W(h) = \lr \left[ (1 - h)\, \wmiss + h\, \whit \right],
\label{eq:writerate}
\end{equation}
and the hit rate is set by whether entries survive their reuse gaps:
\begin{equation}
h^* = F\!\left( S / W(h^*) \right).
\label{eq:fixedpoint}
\end{equation}
Equation~\eqref{eq:fixedpoint} is the classical characteristic-time fixed point with
one addition outside its assumptions: the write rate depends on the
hit rate it produces. The coupling strength is the write
amplification
\begin{equation}
A = \wmiss / \whit = 22.6:
\label{eq:amp}
\end{equation}
one miss evicts as much cache as twenty-two hits. Because $W$ falls as
$h$ rises, the map $\Phi(h) = F(S / W(h))$ is increasing, and an
increasing map can cross the identity more than once: a warm fixed
point (high $h$, low write pressure, long retention) and a cold one
(miss writes flooding the pool, retention too short to hit) can
coexist at the same offered load, separated by an unstable crossing.
Coexistence appears and disappears at the tangency (saddle-node)
condition $\Phi(h) = h$, $\Phi'(h) = 1$. For a two-point reuse-gap mixture
in which a fraction $x$ of requests return after gaps longer than $y$
and the rest return quickly, the warm point $h = 1 - x$ exists iff the
pool covers the long gap at the warm write rate,
\begin{equation}
S \geq \lr\, y \left[ x\, \wmiss + (1 - x)\, \whit \right],
\label{eq:warmexists}
\end{equation}
which inverts to a closed-form collapse edge in load and a matching
recovery edge with the write mix evaluated on the cold branch.
Whether retention feedback alone produces coexistence depends on
where $S / W(0)$ lands in $F$: it must land inside $F$'s central mass. In
a capacity-starved production configuration we analyzed before this
campaign (8 replicas, 0.95M-token pools, 130k-token mean contexts,
$\sim$10 requests/s), $S / W(0) = 11.6$\,s against think times of tens of
seconds, and \eqref{eq:writerate}--\eqref{eq:warmexists} alone give a coexistence band. On the measured
stack of this paper the pools are larger and the protocol caps reuse
gaps at 10.5\,s, so $S / W(0) = 43$--$80$\,s at the studied loads sits far
above every gap: equation~\eqref{eq:fixedpoint} then has a single warm solution at
every studied rate, and no static retention argument can produce the
cold state that Section~\ref{sec:collapse} measures. That state requires a second
feedback.

\subsection{The occupancy-queueing feedback and the cold branch}
\label{sec:occupancy}

Two mechanisms tie the cache to the request queue. First, running
requests occupy the pool: an admitted request reserves its KV need
until completion, and admission fills the pool to a watermark $\theta$
(0.92 in the measured engine) whenever work is waiting. The cache
lives in the remainder, so under a backlog the retention window is
\begin{equation}
\Teff = S\, (1 - \mathrm{occ}) / W(h), \quad \mathrm{occ} \to \theta,
\label{eq:teff}
\end{equation}
a 12.5$\times$ contraction at the watermark. Second, a queued request's
prefix must survive its reuse gap plus its waiting time $w$, so the
survival probability is $F(\Teff - w)$. Both corrections vanish in
light traffic and bite together under a backlog. The service side
closes the loop: with the prefill engine as bottleneck, the fleet
serves at
\begin{equation}
\mu(h) = N p \,/\, \mathrm{E}[\text{prefill per request at } h],
\label{eq:mu}
\end{equation}
which spans a factor of $\sim$27 between branches: $\mu(\text{warm}) \sim 31$
requests/s versus $\mu(0) = N p / \cmiss = 1.16$ requests/s. The
measured cold-state throughput is 1.2--1.4 requests/s \evid{R1}\evid{C01},
within $\sim$15\% of the $\mu(0)$ estimate (the small excess is the measured
2--3\% residual hit rate). A backlog therefore sustains itself
whenever flux into the fleet exceeds $\mu(0)$ while prefixes are cold:
misses hold the pool pinned and the waits long, which holds the hit
rate near zero, which holds the service rate at $\mu(0)$.

Figure~\ref{fig:fixedpoint}a shows the structure as a one-step hit-rate map conditioned
on backlog depth $Q$ at the reference operating point ($N = 8$,
$\ls = 0.022$): with $Q = 0$ the map is monostable warm (the static
result of \ref{sec:retention}); at moderate $Q$ it folds and two stable crossings
coexist; deeper backlogs push the escape threshold - the hit rate a
fleet must already hold for its cache to outrun its queue - toward
one. Figure~\ref{fig:fixedpoint}b integrates the joint $(h, Q)$ dynamics through the
post-perturbation drain window and colors the basins of attraction.
The basin boundary is the separatrix of the drain-versus-rewarm
race; the drain-depth observable of Section~\ref{sec:separatrix} is its empirical
coordinate, KV occupancy read just after the perturbation ends
telling which side of the boundary the trajectory is on.

\begin{figure*}[t]
\centering
\includegraphics[width=\textwidth]{fig_fixed_point.png}
\caption{Theory constructions at ($N{=}8$, $\ls{=}0.022$), workload
functionals and constants from the Section~\ref{sec:calibration} calibration; model
output, no measured data. (a) The one-step hit-rate map $\Phi(u \mid Q)$
conditioned on backlog depth: monostable warm at $Q = 0$, folded
(bistable) at $Q = 60$, escape threshold near one at $Q = 150$. (b)
Basins of attraction of the joint (hit rate, backlog) dynamics
during the post-perturbation drain window; the basin boundary is the
separatrix that the drain-depth observable of Section~\ref{sec:separatrix} estimates.
Time normalization is qualitative (fixed hit-rate relaxation
constant); the figure shows equilibrium and basin structure, not
calibrated transit times.}
\label{fig:fixedpoint}
\end{figure*}

The demand side determines whether the cold branch is reachable.
Requests within a session are completion-coupled - a session issues
its next turn only after the previous returns - so in-flight demand
defers under slowdown instead of accumulating \evid{R7}\evid{C33}, and a
closed population self-stabilizes (Section~\ref{sec:closedloop}). What makes the cold
branch reachable is a demand reservoir that does not yield: sessions
arriving as an exogenous stream park their next turns while the
fleet is slow and release them as catch-up flux when service
resumes. Agentic workloads realize exactly this structure, which is
why the open-loop session-arrival protocol of Section~\ref{sec:method} - and not a
classical closed-loop benchmark - exposes the failure.

\subsection{Break-point estimates}
\label{sec:breakpoints}

\paragraph{Cold-branch feasibility (necessary condition).} A sustained cold
state requires deferred flux above the cold service rate,
$\lr > \mu(0)$:
\begin{equation}
\ls > N p \,/\, (\cmiss\, \Ereq) = 0.0098 \text{ sessions/s}.
\label{eq:coldbound}
\end{equation}
Below bound~\eqref{eq:coldbound} the backlog drains even with every request missing
and recovery is unconditional. The measured protocol recovers 2/2 at
0.012 and 4/4 at 0.017 \evid{C02}, both above the bound - consistent with
\eqref{eq:coldbound} being necessary rather than sufficient: between the bound and
the measured entrapment edge at 0.020--0.022, the reservoir drains
partly serially (completion-coupled), so the effective flux against
$\mu(0)$ sits below $\lr$ and the race of Figure~\ref{fig:fixedpoint}b is winnable. The
gap between bound~\eqref{eq:coldbound} and the measured edge is the workload's
elasticity margin.

\paragraph{Warm-branch feasibility.} The warm branch persists while warm-state
prefill demand fits, which holds to roughly ten times the studied
loads (prefill utilization 0.08 at $\ls = 0.022$). The warm
branch near the band is therefore not destroyed by load; it is
abandoned. Entry at the studied rates is fluctuation-driven - an
eviction-scale perturbation (Section~\ref{sec:protocol}), or the slow occupancy
growth of deepening sessions, carries the state across the
separatrix - which is why outcomes at fixed operating points are
frequencies rather than deterministic responses \evid{C02}\evid{C05}, and why
the coexistence region manifests as hysteresis rather than as a
capacity wall.

\paragraph{Perturbation scale.} A surge writing at rate $W_s$ evicts parked
prefixes on the pool-traversal timescale $S (1 - \mathrm{occ}) / W_s$ - under
two minutes at the measured saturated write rate - so both protocol
surge durations (900\,s and 2700\,s) evict the pool several times
over. What the longer surge changes is the reservoir: deferred
sessions accumulate six times longer and the post-cancellation race
begins from deeper $Q$. The measured duration contrast (900\,s
recovers, 2700\,s entraps, same rate) \evid{C03} is a reservoir-depth
effect, not an eviction-completeness effect, consistent with the
basin geometry of Figure~\ref{fig:fixedpoint}b.

\paragraph{Asymmetry of the transitions.} Collapse is fast because eviction runs
on the pool-overwrite timescale (minutes under a surge); recovery is
slow because cache coverage grows no faster than real time - content
only becomes $T$ seconds old after $T$ seconds - while every miss along
the way regenerates eviction pressure. The hysteresis of Section~\ref{sec:horizon}
is this asymmetry read at steady state.

\subsection{The host tier: capacity sets the regime, restore bandwidth its throughput}
\label{sec:tier-theory}

A host-memory tier of $\Ccpu$ tokens per replica extends retention.
The tier ingests only new writes (restored blocks already reside
there; measured, Section~\ref{sec:congestion}), so its window adds to \eqref{eq:teff}:
\begin{equation}
\Ttot = \Teff + N \Ccpu / \Wnew,
\label{eq:ttot}
\end{equation}
with $\Wnew$ the new-write rate. The naive sizing rule - delete the
cold branch by making $\Ttot(h{=}0)$ exceed the reuse gaps - is
necessary but not the operative constraint, because a tiered fleet
under overload never reaches $h = 0$: it settles into the congested
regime of Section~\ref{sec:tier}, serving hits through the restore channel. Two
separate quantities govern that regime.

\paragraph{Tier capacity sets persistence.} Under congestion the parked working
set recirculates: each live session is touched once per
recirculation time $\trec \sim L / \muc$, with $L$ the live-session count
and $\muc$ the congested completion rate. A session's tier entry
survives to its next touch iff the tier window exceeds $\trec$, giving
the minimum capacity
\begin{equation}
\Ccpu^* \sim \Wnew\, \trec / N.
\label{eq:ccpustar}
\end{equation}
At the measured congested point (tier ingest $\Wnew = 50$--78k
tokens/s \evid{R4}; $\muc = 2.5$--4.3 requests/s over a live population near
200 sessions, so $\trec \sim 50$--90\,s), rule~\eqref{eq:ccpustar} gives 0.4--0.7M tokens
per replica - bracketing the measured persistence boundary: the
0.42M-effective tier loses its external hits and falls through to
cold collapse, while 0.63M and 0.84M persist \evid{C13}. Below $\Ccpu^*$ the
tier only delays collapse; above it, the tier converts collapse into
congestion.

\paragraph{Restore bandwidth sets congested throughput.} Every completion in the
congested regime re-onboards most of a context through the restore
channel, bounding completions at
\begin{equation}
\murestore = N \Br \,/\, (\hext\, \cmiss) \sim 2.2 \text{ requests/s}
\label{eq:murestore}
\end{equation}
at the measured tier-hit share $\hext \sim 0.7$ and $\Br = 23$k tokens/s,
against demand of 2.6 requests/s at $\ls = 0.022$. The regime
therefore runs with its restore channel pinned (measured: 0.81--1.04
of $N \Br$ \evid{C10}) and its queue diverging slowly (measured:
$\sim$0.8/min) - the discriminators of Section~\ref{sec:congestion}. $\Br$ is the measured
congested-state channel throughput, not a hardware ceiling: the two
measured complete recoveries sustained 1.28--1.51$\times$ this rate during
escape \evid{C11}, a distinction that matters for the model failure of
Section~\ref{sec:failures}.

The design consequence is a trade, not a fix: a tier above $\Ccpu^*$
deletes the cold absorbing state (a capacity cliff) and in exchange
the overloaded fleet occupies a congested state whose throughput is
set by restore bandwidth (a bandwidth cliff). Sizing the tier by \eqref{eq:ccpustar}
without provisioning $\Br$ converts one failure mode into another -
the central operational finding of Section~\ref{sec:tier}.

\subsection{Dimensionless structure}
\label{sec:dimensionless}

Five ratios organize the regimes; values at the reference operating
point ($N = 8$, $\ls = 0.022$, untiered unless stated):

\begin{center}
\begin{tabular}{@{}lll@{}}
\toprule
$A$ & $\wmiss / \whit$ & 22.6 \\
$\Lambda$ & $\lr \cmiss / (N p)$ & 2.2 \\
$G$ & $S (1 - \theta) / (W(0)\, g)$ & $\sim$0.4 \\
$\kappa$ & $N \Ccpu / (\Wnew \trec)$ & 0.7--2.0 (sizes 250--500) \\
$\rho_r$ & $\lr \hext \cmiss / (N \Br)$ & $\sim$1.2 \\
\bottomrule
\end{tabular}
\end{center}

$A$ is the feedback strength: $A \sim 1$ is a classical cache; $A \gg 1$ makes
the write rate a state variable. $\Lambda$ is the cold-branch load
ratio: $\Lambda < 1$ forbids a sustained cold state (bound~\eqref{eq:coldbound});
$\Lambda > 1$ makes it feasible and the outcome race-decided. $G$
compares the pinned-pool retention window to the reuse-gap scale $g$:
$G < 1$ means a backlogged fleet cannot hit, which arms the cold
branch; the same fleet at $h = 0$ with its backlog drained has
$G \sim 1.7$, which is why a cold but idle fleet re-warms. $\kappa$ is the
tier-persistence ratio of rule~\eqref{eq:ccpustar}: the measured fall-through sits
at $\kappa < 1$, the persistent congested regime at $\kappa > 1$. $\rho_r$ is
the restore-channel load: $\rho_r \geq 1$ pins the channel and the
congested queue diverges. The router-span dependence of Section~\ref{sec:span} is
deliberately absent from these ratios - it is a property of how one
balancer couples $N$ replicas' drain races (Section~\ref{sec:mechanism}) - and the
theory's scope is the existence, location, and scaling of the
regimes at fixed span.

\subsection{What the theory does not claim}
\label{sec:theory-scope}

The constructions are mean-field and first-order. They do not
predict collapse probability inside the band (outcomes there are
fluctuation-decided; the calibrated simulation of Section~\ref{sec:models} carries
those claims, with documented biases), the magnitude of the span
shift (Section~\ref{sec:models} reproduces its direction emergently), the timing of
delayed collapses, or the escape dynamics of the congested regime
(the model's documented recovery-branch failure, Section~\ref{sec:failures}). Each
formula's anchor - $\mu(0)$ against the measured cold throughput, rule
\eqref{eq:ccpustar} against the measured tier-size boundary, bound~\eqref{eq:murestore} against the
measured restore pin - is a consistency check at one operating
point, not a validated response surface.
